%---------------------------------------------------------------------
\section{Finger Injury Management}
\label{sec:injury}

\begin{warningbox}{Evidence Caveat: No Climbing-Specific Rehabilitation RCTs}
No randomized controlled trials (RCTs) have been conducted on climbing finger injury rehabilitation. Every rehabilitation protocol in this section is classified \textbf{Bioplausible} (mechanistically extrapolated from non-climbing tendinopathy RCTs or tissue biology) or \textbf{Folklore} (practitioner/community heuristic). Evidence-backed claims are restricted to diagnostic imaging accuracy, biomechanics studies, and the Sjöman et al.\ (2023) injury-risk survey. This limitation must be front-loaded in any clinical application of these protocols.
\end{warningbox}

%---------------------------------------------------------------------
\subsection*{Summary Table}

\begin{longtable}{>{\raggedright\arraybackslash}p{3.8cm}
                  >{\raggedright\arraybackslash}p{2.8cm}
                  >{\raggedright\arraybackslash}p{2.2cm}
                  >{\raggedright\arraybackslash}p{3.5cm}
                  >{\raggedright\arraybackslash}p{3.2cm}
                  >{\raggedright\arraybackslash}p{2.2cm}}
\toprule
\textbf{Study} & \textbf{Design / n} & \textbf{Injury type} & \textbf{Protocol / Measure} & \textbf{Primary outcome} & \textbf{Quality} \\
\midrule
\endhead
\bottomrule
\endfoot

Sjöman, Grønhaug \& Julin (2023). \textit{Wilderness \& Environmental Medicine}, 34(4), 397--404. DOI: 10.1016/j.wem.2023.06.004
& Cross-sectional survey; $n=434$ Finnish climbers
& Finger injury (SPIIF onset)
& Association of regular finger training (RFT) with sport-impairing injury by experience level
& RFT negatively assoc.\ with SPIIF in $\geq$6\,yr climbers ($\chi^2=4.57$, $p=0.03$); positively assoc.\ in $<$6\,yr climbers reaching $\geq$7a ($\chi^2=4.61$, $p=0.03$)
& Moderate; self-report; cross-sectional; no causal inference \\

\addlinespace

Klauser et al.\ (2002). \textit{Radiology}. DOI: 10.1148/radiol.2223010752
& Diagnostic accuracy study; $n=64$ climbers, 75 symptomatic fingers
& Flexor pulley injuries (A2/A3/A4)
& Dynamic US (tendon--phalanx distance $>$1.0\,mm during forced flexion) vs.\ MRI reference; surgical correlation in subset
& US sensitivity 98\%, specificity 100\% vs.\ MRI; verification bias acknowledged
& Moderate; reference standard non-uniform \\

\addlinespace

Berrigan et al.\ (2021). \textit{Journal of Ultrasound Medicine}. DOI: 10.1002/jum.15796
& Narrative/systematic imaging review
& A2 pulley diagnosis
& US vs.\ MRI/CT; dynamic maneuvers
& US recommended as primary modality; sensitivity/specificity 0.94--1.00 across included studies
& Moderate; secondary synthesis \\

\addlinespace

Schöffl VR et al.\ (2003). \textit{Wilderness \& Environmental Medicine}, 14(2), 94--100.
& Prospective clinical registry cohort; $n=604$ climbers, 122 pulley injuries
& Pulley injuries Grade I--IV
& Conservative (Grades I--III: immobilization $\pm$ tape $\pm$ ring) vs.\ surgical (Grade IV loop reconstruction)
& 92\% conservative Grade I--III returned to prior level; Grade IV surgical outcomes largely good; no controls
& Low--moderate; case series; no RCT \\

\addlinespace

Schweizer A.\ (2000). \textit{Journal of Hand Surgery (Br)}, 25(3), 243--246. DOI: 10.1054/jhsb.1999.0335
& Biomechanical in-vivo; $n=16$ fingers
& A2 pulley loading / taping
& Circular taping variants in crimp grip; tendon--bone distance measurement
& Tape over proximal end A2: bowstringing $\downarrow$2.8\%, force absorbed $\approx$11\%; more distal placement: bowstringing $\downarrow$22\%, force absorbed $\approx$12\%
& Moderate for biomechanics; no clinical outcome data \\

\addlinespace

Alfredson H et al.\ (1998). \textit{American Journal of Sports Medicine}, 26(3), 360--366. DOI: 10.1177/03635465980260030301
& Prospective cohort + historical controls; $n=15+15$
& Chronic Achilles tendinosis (tendon analogue)
& Heavy-load eccentric heel-drop $3{\times}15$ reps, twice daily, 12\,wk vs.\ continued activity
& 100\% return to running at 12\,wk (eccentric) vs.\ 0\% (control at study midpoint); pain and tendon structure normalised
& High for Achilles tendinopathy; extrapolation to pulleys bioplausible only \\

\addlinespace

Cook JL \& Purdam CR (2009). \textit{British Journal of Sports Medicine}, 43(6), 409--416. DOI: 10.1136/bjsm.2008.051193
& Conceptual review / mechanistic model
& Tendinopathy (general)
& Reactive $\to$ dysrepair $\to$ degenerative continuum; load-management framework
& Reactive stage: load reduction not rest; dysrepair/degenerative: progressive isotonic loading
& Moderate; well-cited model; not RCT \\

\addlinespace

Rio E et al.\ (2015). \textit{British Journal of Sports Medicine}, 49(19), 1277--1283. DOI: 10.1136/bjsports-2014-094386
& RCT crossover; $n=6$ per condition (patellar tendinopathy)
& Patellar tendinopathy (analogue)
& Single-bout isometric vs.\ isotonic leg extension
& Isometric: pain $\downarrow$6.8/10 immediately (7.0$\pm$2.04 $\to$ 0.17$\pm$0.41), sustained 45\,min; MVIC $\uparrow$18.7$\pm$7.8\%
& Moderate; very small $n$; single session; extrapolation to pulleys bioplausible only \\

\addlinespace

Shaw G et al.\ (2017). \textit{American Journal of Clinical Nutrition}, 105(1), 136--143. DOI: 10.3945/ajcn.116.138594
& RCT crossover; $n=8$
& Connective tissue (ligament model)
& 15\,g gelatin + 48\,mg vitamin\,C, 1\,hr pre-exercise vs.\ placebo
& P1NP (collagen synthesis marker) doubled vs.\ placebo; in-vitro engineered ligament mechanics improved
& Moderate; very small $n$; surrogate endpoints only; no injury outcome data \\

\addlinespace

Iruretagoiena et al.\ (2020). \textit{Wilderness \& Environmental Medicine}. DOI: 10.1016/j.wem.2020.07.007
& Systematic review; 7 studies
& A2/A4 pulley diagnosis
& Ultrasound tendon--bone distance methods
& No consensus thresholds; complete A2 rupture cutoffs range 1.9--5.1\,mm across studies
& Moderate diagnostic review \\

\end{longtable}

%---------------------------------------------------------------------
\subsection{Injury-Risk Stratification Framework}
\label{subsec:riskstrata}

Risk tiers used throughout this section. Incidence estimates derive from available survey and case series data; where climbing-specific data are absent, estimates are flagged as extrapolated from adjacent sports.

\renewcommand{\arraystretch}{1.3}
\begin{center}
\small
\begin{tabular}{>{\raggedright}p{2.4cm}
                >{\raggedright}p{3.0cm}
                >{\raggedright}p{4.0cm}
                >{\raggedright\arraybackslash}p{5.0cm}}
\toprule
\textbf{Stratum} & \textbf{Estimated incidence} & \textbf{Definition} & \textbf{Management implication} \\
\midrule
\textbf{Low} & $<$5\,\%/year & Managed with standard warm-up and load monitoring; no elevated structural risk & Routine load progression; standard 48\,h recovery; no additional screening \\
\addlinespace
\textbf{Moderate} & 5--20\,\%/year & Context-dependent risk; structural risk present but mitigable with specific precautions & Explicit monitoring of finger pain VAS; HRV supplementary tracking; grip-position modification if VAS $\geq$2/10 \\
\addlinespace
\textbf{High} & $>$20\,\%/year & Elevated structural risk; explicit mitigation required & Mandatory load reduction, grip-position restriction, structured rehabilitation; medical evaluation if VAS $\geq$3/10 at rest \\
\addlinespace
\textbf{Very High} & Near-certain injury without intervention & Contraindicated without medical supervision & Full climbing rest pending medical evaluation; return protocol requires clinical clearance \\
\bottomrule
\end{tabular}
\end{center}

\noindent\textbf{Note on incidence data.} No prospective climbing study has measured annual pulley injury incidence by stratum. The rate estimates above are informed by Schöffl et al.\ (2003) and Sjöman et al.\ (2023) survey data and should be treated as indicative. Injury incidence and return-to-sport risk are separate quantities: a Low-incidence injury can carry High return-to-sport risk if mismanaged. These strata are stated explicitly at each recommendation below where applicable.

%---------------------------------------------------------------------
\subsection{A2 Pulley Injuries}
\label{subsec:a2pulley}

\subsubsection{Anatomy and Mechanism}

The flexor tendon sheath of digits II--V contains five annular (A1--A5) and three cruciate (C1--C3) pulleys. The A2 pulley spans the proximal half of the proximal phalanx and is the largest, least deformable annular pulley; it prevents bowstringing of the flexor digitorum profundus (FDP) and flexor digitorum superficialis (FDS) by maintaining the tendons' mechanical line of action. In published case series, the ring (digit IV) and middle (digit III) fingers account for $\approx$98\% of pulley injuries in climbers, consistent with their role as the primary force-bearing digits in crimp configurations.

\begin{bioplausbox}{Full-Crimp Loading Hierarchy: A2 Pulley Contact Pressure · BIOPLAUSIBLE · ESS 2}
\textbf{Grip-position loading hierarchy (Bioplausible):} Full crimp (PIP $\approx$90$^\circ$ flexion, DIP hyperextension or 0--20$^\circ$) $>$ half crimp $>$ open hand, due to acute tendon approach angle against the pulley increasing contact pressure. Schweizer (2001) estimated full-crimp A2 forces at 35--40\,N/kg body mass in cadaveric/model studies, half-crimp at $\approx$30--35\,N/kg, and open-hand at $\approx$20--25\,N/kg --- a $\approx$30--50\% load reduction comparing full crimp to open hand. No 95\,\% CI was reported. These data underpin the common recommendation to favour open-hand grip during volume accumulation; however, no prospective injury-prevention RCT has confirmed that grip-position prescription reduces pulley injury incidence in climbers. \textbf{[Bioplausible: mechanistic from Schweizer 2000 cadaveric biomechanics; not confirmed in prospective trial.]}
\end{bioplausbox}

\subsubsection{Grading}

The Schöffl (2003) grading system remains the dominant clinical and research classification:

\begin{itemize}[leftmargin=1.5em]
  \item \textbf{Grade I}: Partial tear; no bowstringing on static or dynamic US. Tenderness over A2; pain with resisted PIP flexion.
  \item \textbf{Grade II}: Partial tear with mild bowstringing ($<$2\,mm tendon-to-bone distance [TBD] increase under load on US). Palpable swelling; pain with active loading.
  \item \textbf{Grade III}: Complete A2 tear; bowstringing clearly visible ($>$2\,mm TBD increase). May have bruising and palpable defect.
  \item \textbf{Grade IV}: Complete A2 tear plus adjacent pulley involvement (typically A3 or A4) or FDP/FDS damage. Greater bowstringing; potential flexor tendon dysfunction.
\end{itemize}

\begin{bioplausbox}{Schöffl Grading System: Clinical Utility Without RCT Validation · BIOPLAUSIBLE · ESS 2}
\textbf{Grading validity (Bioplausible):} The Schöffl system derives from case series, not a validated prospective grading-outcome study. No inter-rater reliability data from RCTs exists. The system is clinically useful for mapping structural severity to management pathway but should not be treated as a validated prognostic instrument. \textbf{[Bioplausible: supported by case series, Schöffl et al.\ 2003; no RCT validation.]}
\end{bioplausbox}

\subsubsection{Diagnosis}

\paragraph{Physical examination.}
Pain reproduced by resisted PIP flexion and direct palpation over the proximal volar aspect of the proximal phalanx is the primary clinical indicator. Bowstringing may be palpable or visible in Grade III--IV when the digit is actively flexed against resistance. Sensitivity and specificity of clinical examination alone are not established in controlled diagnostic accuracy studies.

\begin{bioplausbox}{Clinical Examination Reliability for Pulley Injury: Unvalidated · BIOPLAUSIBLE · ESS 2}
\textbf{Clinical exam reliability (Bioplausible):} Focal tenderness and pain with resisted PIP flexion are widely used but unvalidated against imaging or surgical reference standards in controlled studies. \textbf{[Bioplausible: clinical reasoning; no ROC/diagnostic accuracy study.]}
\end{bioplausbox}

\paragraph{Imaging.}

\begin{evidencebox}{Ultrasound as Primary Diagnostic Modality for Pulley Injury · EVIDENCE\_BACKED · ESS 6}
\textbf{Ultrasound as primary modality (Evidence-backed):} Dynamic high-resolution US ($\geq$15\,MHz linear probe, assessed in passive extension and full resisted flexion) is the recommended first-line modality. Klauser et al.\ (2002) reported US sensitivity 98\%, specificity 100\% vs.\ MRI reference standard in 64 climbers with 75 symptomatic fingers (DOI: 10.1148/radiol.2223010752; verification bias acknowledged). Berrigan et al.\ (2021) systematic review confirms US sensitivity/specificity 0.94--1.00 across included studies (DOI: 10.1002/jum.15796). Dynamic assessment comparing TBD at 20$^\circ$ passive extension vs.\ full active flexion is essential; static views underestimate Grade II bowstringing and miss partial tears.

Complete A2 rupture TBD diagnostic cutoffs are inconsistent across studies, ranging 1.9--5.1\,mm (Iruretagoiena et al.\ 2020, DOI: 10.1016/j.wem.2020.07.007). MRI sensitivity for complete pulley tears is approximately 90--95\% and for partial tears 65--80\% in the general hand surgery literature, but these figures are not from climbing-specific populations. MRI is indicated when US is equivocal or multi-structure injury is suspected. \textbf{Bottomline: use high-resolution dynamic US first; reserve MRI for equivocal or complex Grade IV presentations.}
\end{evidencebox}

%---------------------------------------------------------------------
\subsubsection{Conservative Management --- Grade I--II and Most Grade III}

\begin{warningbox}[title=No Climbing Rehabilitation RCTs]
All immobilisation durations, return-to-loading timelines, and hangboard progression thresholds below are derived from case series (Schöffl et al.\ 2003) or extrapolated from non-climbing tendinopathy RCTs. None has been validated in a controlled climbing study.
\end{warningbox}

\paragraph{Immobilisation.}\leavevmode

\begin{bioplausbox}{Short-Term Immobilisation: Grade-Based Duration (Schöffl 2003) · BIOPLAUSIBLE · ESS 2}
\textbf{Short-term immobilisation (Bioplausible):} Typical recommendations from Schöffl et al.\ (2003) case series:
\begin{itemize}[leftmargin=1.5em,noitemsep]
  \item Grade I: 0\,d immobilisation; functional therapy 2--4\,wk; tape protection.
  \item Grade II: 10\,d immobilisation; then 2--4\,wk functional therapy.
  \item Grade III: 14\,d immobilisation with thermoplastic ring; then up to 4\,wk functional therapy.
\end{itemize}
Rationale: early controlled mechanical stimulation promotes aligned collagen deposition; prolonged immobilisation ($>$1\,wk) is inconsistent with tendon and ligament healing biology. However, no climbing RCT has compared immobilisation durations. \textbf{[Bioplausible: extrapolated from tendon biology; case series source.]}
\end{bioplausbox}

\paragraph{Return-to-loading timeline.}\leavevmode

\begin{bioplausbox}{Return-to-Climbing Timelines After Pulley Injury (Schöffl 2003 Case Series) · BIOPLAUSIBLE · ESS 2}
\textbf{Return-to-climbing timelines (Bioplausible):} Reported in Schöffl et al.\ (2003) case series; not validated in controlled trials.
\begin{itemize}[leftmargin=1.5em,noitemsep]
  \item Grade I: progressive reloading at 1--2\,wk; return to climbing 4--6\,wk.
  \item Grade II: progressive reloading at 3--4\,wk; return to climbing 6--10\,wk.
  \item Grade III (conservative): progressive reloading at 6--8\,wk; return to full climbing 3--5\,months.
\end{itemize}
\textbf{[Bioplausible: mechanistically consistent with tendon healing biology; Folklore for specific week numbers (Schöffl 2003 case series).]}
\end{bioplausbox}

\paragraph{H-taping and ring splinting.}\leavevmode

\begin{bioplausbox}{H-Taping: Biomechanical Bowstringing Reduction · BIOPLAUSIBLE · ESS 2}
\textbf{Biomechanical basis (Bioplausible):} Schweizer (2000) in-vivo study ($n=16$ fingers) found tape placed directly over the A2 pulley reduced bowstringing by 2.8\% and absorbed $\approx$11\% of bowstringing force; more distal placement (toward distal end of proximal phalanx) reduced bowstringing by 22\% and absorbed $\approx$12\%. Schöffl et al.\ (2007) controlled biomechanical study ($n=8$ US, $n=12$ strength assessment)\footnote{Single-source claim --- requires independent verification.} found H-tape reduced tendon-bone distance by 16\% and increased crimp MVC by $+$13\% (no CI reported) vs.\ no tape; no effect in open-hand hang. A systematic review by Larsson et al.\ (2022) reported low-to-moderate certainty for bowstringing reduction (15--22\%) in injured fingers; no significant effect in uninjured fingers. \textbf{[Bioplausible: small controlled biomechanics data; no RCT for clinical outcomes.]}
\end{bioplausbox}

\begin{folklorebox}{H-Taping in Climbing Rehabilitation: Universal Practice · FOLKLORE\_VERIFIED · ESS 1}
\textbf{H-taping clinical practice (Folklore):} H-taping around the proximal phalanx (figure-8 configuration creating dorsal bridge) is near-universal in climbing rehabilitation practice. Community consensus, Hooper's Beta, Schöffl-era clinical guidelines. No RCT or controlled case series has measured re-injury rates, healing time, or functional outcomes with vs.\ without taping. \textbf{[Folklore: climbing community / climbing physio consensus.]}
\end{folklorebox}

\paragraph{Load-managed rehabilitation --- tendinopathy extrapolation.}\leavevmode

\begin{evidencebox}{Tendon Loading Paradigms: Eccentric and Isometric Evidence (Achilles/Patellar) · EVIDENCE\_BACKED · ESS 3}
\textbf{Tendon loading paradigms (Evidence-backed for analogous tissues):}
\begin{itemize}[leftmargin=1.5em]
  \item Alfredson et al.\ (1998, DOI: 10.1177/03635465980260030301): heavy-load eccentric heel-drop (3$\times$15 twice daily, 12\,wk) achieved 100\% return to running in $n=15$ chronic Achilles tendinosis patients vs.\ 0\% for continued activity alone. A 2023 meta-analysis of 5 Achilles RCTs reported eccentric exercise pain reduction MD $-$1.21 (95\,\% CI $-$2.72 to $-$0.30) vs.\ other exercise.
  \item Cook \& Purdam (2009, DOI: 10.1136/bjsm.2008.051193): reactive--dysrepair--degenerative continuum model; reactive stage warrants load reduction (not rest) and isometric loading; degenerative stage benefits from progressive isotonic loading.
  \item Rio et al.\ (2015, DOI: 10.1136/bjsports-2014-094386): isometric leg extension produced immediate pain reduction of 6.8/10 (from 7.0$\pm$2.04 to 0.17$\pm$0.41, sustained 45\,min) and MVIC $\uparrow$18.7$\pm$7.8\% vs.\ isotonic in patellar tendinopathy ($n=6$).
\end{itemize}
\textbf{[Evidence-backed for Achilles and patellar tendinopathy. Extrapolation to A2 pulley is Bioplausible only --- see below.]}
\end{evidencebox}

\begin{bioplausbox}{Progressive Loading Applied to A2 Pulley: Extrapolation from Tendinopathy · BIOPLAUSIBLE · ESS 2}
\textbf{Pulley-specific loading (Bioplausible):} The A2 pulley is a fibrocartilaginous ligamentous condensation of the flexor sheath --- primarily type\,I collagen with fibrocartilaginous zones under compressive load --- not a tendon. Cellular overlap with tendons exists (fibroblast mechanosensing, integrin-mediated collagen remodelling), but load-response properties are not identical. Application of progressive loading protocols (isometrics $\to$ isotonic $\to$ functional) to pulley rehabilitation is bioplausible but constitutes an extrapolation from a different tissue type. No peer-reviewed study has measured A2 pulley load during hangboard exercises with validated force-sensing systems and correlated measurements with tissue healing outcomes. \textbf{[Bioplausible: mechanistic framework from Cook \& Purdam 2009 and Alfredson 1998; Folklore: specific thresholds and protocols.]}
\end{bioplausbox}

\paragraph{Load-managed hangboard rehabilitation during recovery.}\leavevmode

\begin{folklorebox}{Tyler Nelson Progressive Loading Protocols for Pulley Rehabilitation · FOLKLORE\_VERIFIED · ESS 1}
\textbf{Tyler Nelson / Camp 4 Human Performance protocols (Folklore):} Tyler Nelson has disseminated progressive loading protocols for pulley rehabilitation through clinical practice and online/professional channels. Proposed structure: (a) isometric hangs beginning at 20--40\% perceived maximum in the early reactive phase, pain $\leq$2--3/10; (b) progress load 5--10\% per week guided by symptom response; (c) grip-position progression --- open-hand $\to$ half crimp, full crimp avoided until $\approx$week 12. These thresholds are not empirically derived from pulley-specific loading studies. \textbf{[Folklore: Tyler Nelson, C4HP, practitioner/coaching channels. Mechanistic framework Bioplausible via Cook \& Purdam 2009; specific thresholds unvalidated.]}
\end{folklorebox}

%---------------------------------------------------------------------
\subsubsection{Surgical Management --- Grade III--IV}

\begin{bioplausbox}{Surgical Indications for Grade III--IV Pulley Injury · BIOPLAUSIBLE · ESS 2}
\textbf{Surgical indications (Bioplausible / case series):} Surgery is indicated for: (a) Grade IV with documented tendon involvement; (b) Grade III in elite/professional athletes requiring accelerated return or where conservative management ($\geq$4--6\,months) has failed; (c) persistent symptomatic bowstringing with functional deficit. Techniques include direct suture repair, loop reconstruction using palmaris longus or extensor retinaculum graft (Lister, Kleinert/Weilby), and loop-and-a-half reconstruction.

Arora et al.\ (2007) retrospective comparative series ($n=23$) reported excellent/good Buck-Gramcko outcomes in 21/23 with no major complications.\footnote{Single-source claim --- requires independent verification.} Bouyer et al.\ (2016) retrospective series ($n=38$, mean FU 85\,months)\footnote{Single-source claim --- requires independent verification.} found post-operative US E-space $<$2\,mm predicted return to prior climbing level in 70\% vs.\ 25\% in those not returning. Re-injury rates are approximately 15--25\% in case series follow-up. Return-to-climbing: 3--6\,months Grade III; 6--12\,months Grade IV; mean 5.8\,months to pre-injury level reported in one elite-climber series.\footnote{Single-source claim --- requires independent verification.}

\textbf{[Bioplausible / Low-quality evidence: all data from case series; no controls; heterogeneous surgical techniques; publication bias likely.]}
\end{bioplausbox}

%---------------------------------------------------------------------
\subsection{Flexor Tendon Strains}
\label{subsec:flexortendon}

\begin{bioplausbox}{Flexor Tendon Strain: Anatomy and Mechanism · BIOPLAUSIBLE · ESS 2}
\textbf{Anatomy and mechanism (Bioplausible):} FDP inserts distally to the distal phalanx; FDS splits at the proximal phalanx to insert bilaterally onto the middle phalanx base. Both run through the same fibrous sheath as the A2 pulley. Flexor tendon strains (partial or microtear injuries to the tendon proper) typically occur at the musculotendinous junction or within the tendon substance under repeated high-load crimp grip loading or sudden eccentric overload. They co-occur with Grade IV pulley injuries (definitional) or may present independently. No climbing-specific epidemiological data distinguish flexor tendon strain prevalence from pulley injury prevalence in isolation. \textbf{[Bioplausible: anatomical reasoning and clinical report.]}
\end{bioplausbox}

\paragraph{Grading and differentiation from pulley injury.}\leavevmode

\begin{bioplausbox}{Flexor Tendon Strain Grading and Differentiation from Pulley Injury · BIOPLAUSIBLE · ESS 2}
\textbf{Grading (Bioplausible / clinical consensus):}
\begin{itemize}[leftmargin=1.5em,noitemsep]
  \item Grade I: minor fibre disruption; no architectural loss on US; tenderness along tendon course; pain with resisted flexion.
  \item Grade II: partial tear; hypoechoic defect on US; weakened active flexion; focal swelling.
  \item Grade III: complete/near-complete disruption; functional deficit; requires urgent surgical referral.
\end{itemize}
Distinguishing from pulley injury: pulley injury presents with focal tenderness over the proximal phalanx with bowstringing on dynamic US; flexor tendon strain presents with pain along the tendon course proximal or distal to the pulley, internal tendon changes on US, but no bowstringing. Precise attribution requires high-quality dynamic US. \textbf{[Bioplausible: no climbing-specific grading validation study.]}
\end{bioplausbox}

\paragraph{Management.}\leavevmode

\begin{bioplausbox}{Grade I--II Flexor Tendon Strain: Conservative Management · BIOPLAUSIBLE · ESS 2}
\textbf{Grade I--II management (Bioplausible):} Principles mirror A2 pulley Grade I--II management: (a) acute load reduction 5--14\,d; (b) early progressive loading from week 2--3, pain-guided; (c) progressive return to full loading over 6--12\,wk. Mechanistic basis is identical to the Cook--Purdam continuum framework. No climbing-specific RCT or case series specifically distinguishes flexor tendon strain outcomes from pulley injury outcomes. \textbf{[Bioplausible: same extrapolation from tendinopathy literature.]}
\end{bioplausbox}

Grade III FDP avulsion (``jersey finger'' equivalent): requires urgent surgical referral and formal hand surgery protocol --- outside scope of conservative climbing rehabilitation.

%---------------------------------------------------------------------
\subsection{Skin Management}
\label{subsec:skin}

\subsubsection{Flappers --- Acute Traumatic Skin Tears}

A flapper is a traumatic partial-thickness skin tear, occurring at the callus-to-normal-skin transition where a stress riser under tangential shear load on rock hold edges causes delamination. Most common on the second and third digit palmar skin and thenar eminence.

\begin{folklorebox}{Acute Flapper Management: Universal Climbing Practice · FOLKLORE\_VERIFIED · ESS 1}
\textbf{Acute management (Folklore):} Trim any non-viable detached skin flap with sterile scissors (reduces bacterial trapping and maceration); irrigate; apply occlusive dressing or thin skin balm (petrolatum, Joshua Tree Salve, Climb On, or similar humectant) to maintain moist wound environment. Taping over a partially healed wound allows earlier return. \textbf{[Folklore: climbing community universal practice. General wound healing biology supports moist wound environment for epidermal regeneration (bioplausible extrapolation), but no climbing-specific trial exists.]}
\end{folklorebox}

\begin{toverifybox}{Return-to-Climbing Timeline After Flappers · TO\_VERIFY · ESS 1}
\textbf{TO\_VERIFY:} Anecdotal community consensus: mild flapper (superficial, $<$1\,cm$^2$) 3--7\,d; moderate (deeper or larger) 7--14\,d; larger raw dermis 1--2\,wk --- dependent on pain, infection risk, and hold texture. No empirical healing timeline data exist. Requires $\geq$3 independent named climbing sources with specific timelines for FOLKLORE\_VERIFIED status.
\end{toverifybox}

\begin{folklorebox}{Callus Management for Flapper Prevention · FOLKLORE\_VERIFIED · ESS 1}
\textbf{Prevention --- callus management (Folklore):} Filing calluses flush (reducing height and smoothing the callus-to-normal-skin transition) is near-universal in elite climbing practice; mechanistic goal is to reduce the stress riser at the thickness transition. Has not been evaluated in any study measuring flapper incidence as a function of callus management. Additional community recommendations: moisturise nightly, avoid alcohol-based sanitisers before climbing, manage session volume on sharp holds, rotate grip styles. \textbf{[Folklore: Eric Horst, \textit{Training for Climbing}; Lattice Training / Will Anglin; athlete and coach consensus.]}
\end{folklorebox}

\subsubsection{Chronic Skin Thinning}

\begin{bioplausbox}{Chronic Skin Thinning: Mechanism from Repeated Hold Friction · BIOPLAUSIBLE · ESS 2}
\textbf{Mechanism (Bioplausible):} Repeated friction on small-radius holds reduces stratum corneum thickness below functional baseline, presenting as hypersensitivity, split skin, and reduced friction coefficient. Contributing factors: training volume, hold substrate roughness, ambient humidity (very low humidity accelerates transepidermal water loss and skin cracking), and individual skin phenotype. No climbing-specific study has measured stratum corneum thickness longitudinally as a function of training load. \textbf{[Bioplausible: mechanism from general dermatology; no climbing-specific intervention data.]}
\end{bioplausbox}

\begin{folklorebox}{Climbing Skin Products and Humidity Management · FOLKLORE\_VERIFIED · ESS 1}
\textbf{Interventions (Folklore):} Climbing-specific skin products (Rhino Skin Solutions, ClimbSkin, Friction Labs products), humidity management (40--60\% indoor humidity anecdotally preferred), Antihydral (methenamine; acts by protein denaturation to reduce sweating, potentially increasing friction but risking ``glassy'' skin and splits), callus sanding, and session spacing are community-recommended. No peer-reviewed controlled trial has evaluated any climbing skin product for skin health, injury rate, or return-to-climbing outcomes. \textbf{[Folklore: manufacturer claims, community/practitioner consensus; Dave MacLeod, ClimbSkin. Emollient/humectant mechanisms are Bioplausible from general dermatology literature on skin barrier repair.]}
\end{folklorebox}

\subsubsection{Collagen Supplementation for Skin Healing}

\begin{evidencebox}{Collagen Supplementation: Skin and Connective Tissue Endpoints in Non-Climbing Populations · EVIDENCE\_BACKED · ESS 3}
\textbf{Collagen synthesis markers and skin endpoints (Evidence-backed for non-climbing populations):}
\begin{itemize}[leftmargin=1.5em]
  \item Shaw et al.\ (2017, DOI: 10.3945/ajcn.116.138594): RCT crossover ($n=8$); 15\,g gelatin + 48\,mg vitamin\,C, 1\,hr pre-exercise vs.\ placebo doubled circulating P1NP (collagen synthesis marker); in-vitro engineered ligament mechanics improved. \textit{Note: $n=8$ is very small; endpoints are surrogate; no injury or skin-healing outcome data.}
  \item A meta-analysis of oral hydrolyzed collagen peptides (26 RCTs, $n=1721$) for skin aging found improved hydration and elasticity endpoints. A systematic review/meta-analysis of 19 RCTs ($n=1341$) on oral/topical peptides found improved hydration and modest wrinkle reduction.
  \item Proksch et al.\ (2025): RCT ($n=66$), 2.5\,g collagen peptides vs.\ placebo --- reduced wrinkles, improved elasticity and hydration.
  \item A double-blind pilot RCT in burn patients ($n=30$): hydrolyzed collagen supplementation improved wound healing markers and reduced hospital stay $\approx$30\% vs.\ control.\footnote{Single-source claim --- requires independent verification.}
\end{itemize}
\textbf{[Evidence-backed for skin elasticity/hydration endpoints in aging and burn populations. Bioplausible for climbing skin healing: no direct evidence for flapper healing or friction-induced thinning in climbers.]}
\end{evidencebox}

%---------------------------------------------------------------------
\subsection{General Injury Prevention Principles}
\label{subsec:prevention}

\subsubsection{Warm-Up}

\begin{bioplausbox}{Warm-Up Efficacy for Injury Prevention: Mechanism Without Climbing RCT · BIOPLAUSIBLE · ESS 2}
\textbf{Warm-up efficacy (Bioplausible):} No climbing-specific RCT exists for injury prevention via warm-up. General sports medicine literature supports warm-up via increased connective tissue temperature (improving viscoelastic properties), increased neuromuscular activation, and enhanced motor pattern recruitment. The most methodologically robust proximate-domain evidence is the FIFA 11+ protocol (Soligard et al.\ 2008), which demonstrated $\approx$50\% lower-limb injury reduction in a large cluster-RCT in soccer --- tissue type, injury mechanism, and loading characteristics differ substantially from climbing finger pulley injuries.

Practical climbing warm-up: 5--10\,min aerobic movement; easy traversing; progressive finger recruitment; submaximal hangs before limit bouldering. \textbf{[Bioplausible: mechanistic basis exists; climbing-specific RCT evidence absent.]}
\end{bioplausbox}

\subsubsection{Load Management --- The 10\% Rule}

\begin{toverifybox}{The 10\,\% Load Progression Rule Applied to Climbing · TO\_VERIFY · ESS 1}
\textbf{TO\_VERIFY --- 10\% rule (unvalidated in climbing):} The heuristic that weekly training load increases should not exceed 10\% of the prior week's load is widely cited in running medicine and climbing coaching. Its evidence base is weaker than commonly assumed: it derives primarily from retrospective epidemiological observation, not prospective RCT. Systematic reviews (e.g., Impellizzeri et al.\ 2020) note that the specific 10\% threshold is arbitrary and inconsistently supported across sports. Gabbett (2016, DOI: 10.1136/bjsports-2015-095788) and the acute:chronic workload ratio literature provide a more quantitative framework but remain primarily observational, contested in threshold values, and derived from team sports. For climbing, the threshold is additionally complicated because intensity, grip type, board angle, session density, and concurrent finger-training volume interact in ways not captured by a single linear loading metric. Requires $\geq$3 independent climbing-community sources citing this specific threshold for FOLKLORE\_VERIFIED status.
\end{toverifybox}

\subsubsection{Experience-Dependent Injury Risk}

\begin{evidencebox}{Experience-Dependent Finger Training Risk (Sjöman 2023) · EVIDENCE\_BACKED · ESS 5}
\textbf{Sjöman et al.\ (2023) --- Experience-dependent association (Evidence-backed):} Cross-sectional survey of $n=434$ Finnish climbers (DOI: 10.1016/j.wem.2023.06.004) examined the association between route/bouldering-specific finger strength training (RFT) and sport/performance-impairing finger injury (SPIIF):
\begin{itemize}[leftmargin=1.5em,noitemsep]
  \item In climbers with $\geq$6\,yr experience: RFT \textbf{negatively} associated with SPIIF ($\chi^2=4.57$, $p=0.03$) --- structured training may be protective.
  \item In climbers with $<$6\,yr experience who had reached $\geq$7a (French sport): RFT \textbf{positively} associated with SPIIF ($\chi^2=4.61$, $p=0.03$) --- training in rapidly progressing novices may increase injury risk.
\end{itemize}
\textit{Quality constraints: cross-sectional design prohibits causal inference; self-reported injury and training data; Finnish climbing population; no dose-response data; no adjustment for confounders (training volume, hold type, warm-up, concurrent activities).}
\end{evidencebox}

\begin{bioplausbox}{Bidirectional Injury Risk Pattern: Mechanistic Interpretation · BIOPLAUSIBLE · ESS 2}
\textbf{Mechanistic interpretation (Bioplausible):} The bidirectional pattern is consistent with theoretical models of training-load injury risk in developing athletes (rapid capacity acquisition outpacing structural tissue adaptation) and the protective adaptation hypothesis in experienced athletes (accumulated structural adaptation of the fibrous flexor sheath, pulley hypertrophy, and neuromuscular co-contraction patterns). The mechanism for protection in experienced climbers has not been directly measured. \textbf{[Bioplausible for mechanism; Evidence-backed for the observed association in this dataset only.]}
\end{bioplausbox}

%---------------------------------------------------------------------
\subsection{What the Literature Cannot Yet Answer}
\label{subsec:injury-gaps}

\begin{enumerate}[leftmargin=2em]
  \item \textbf{Progressive loading vs.\ rest-only for Grade I--II A2 injuries.} Does structured hangboard rehabilitation outperform rest for time-to-return, re-injury rate, or ultrasound-confirmed healing in Grade I--II tears? No RCT exists; a moderate-sized trial ($n \geq 60$, Grade I--II, ultrasound-confirmed outcomes) is feasible.

  \item \textbf{Optimal immobilisation duration and type.} For Grade II--III A2 injuries, does 0--7\,d vs.\ 10--14\,d immobilisation (with standardised pulley protection) alter long-term TBD, return timeline, or re-injury rate? No controlled comparison exists.

  \item \textbf{H-taping clinical effectiveness.} Schweizer (2000) established small biomechanical reductions; Schöffl et al.\ (2007) reported acute crimp-strength improvement. Neither confirms reduced re-injury rate or faster healing in a controlled outcome study. Does H-taping reduce bowstringing enough to change clinical outcomes, or does it act primarily via pain modulation and confidence?

  \item \textbf{Ultrasound-guided return-to-climb criteria.} Does using serial TBD measurement to gate progression (e.g., TBD $<$2\,mm under resisted flexion before advancing grip position) reduce re-injury relative to symptom-only progression? The current TBD cutoff range of 1.9--5.1\,mm for complete rupture precludes standardised criteria.

  \item \textbf{Load-to-failure properties and healing kinetics of the human A2 pulley under progressive loading.} No in-vivo or ex-vivo study has quantified how incremental hangboard loading affects pulley tissue mechanics during healing. Tyler Nelson--style protocols lack any quantitative validation in the target tissue.

  \item \textbf{Collagen supplementation for pulley and skin healing.} Shaw et al.\ (2017) demonstrated elevated P1NP surrogate markers ($n=8$). No study has measured actual A2 pulley recovery time, re-injury rate, or imaging-based healing in climbers using gelatin + vitamin\,C. Similarly, no RCT has evaluated oral collagen peptides for flapper healing or chronic skin thinning under climbing-specific friction trauma.

  \item \textbf{Structural mediators of experience-dependent injury risk.} The Sjöman et al.\ (2023) bidirectional association is observational. A prospective cohort tracking novices from onset of structured finger training through 2--5\,yr, with serial ultrasound pulley morphology assessment, grip-force measurement, and load monitoring, would provide the first mechanistic evidence for whether pulley hypertrophy or neuromuscular adaptation mediates the protective effect observed in experienced climbers.
\end{enumerate}
